\subsection{$L^1$ 函数的 Fourier 变换}
\label{sec:ch1-fourier-transform-l1}

给定函数 $f\in L^1(\mathbb R^n)$, 定义其 Fourier 变换为
\begin{equation}
\label{eq:ch1-fourier-transform-definition}
 \widehat f(\xi)=\int_{\mathbb R^n}f(x)e^{-2\pi ix\cdot\xi}\,dx,
\end{equation}
其中 $x\cdot\xi=x_1\xi_1+x_2\xi_2+\cdots+x_n\xi_n$. 下面列出 Fourier 变换的若干性质:
\begin{equation}
\label{eq:ch1-fourier-linearity}
 \widehat{(\alpha f+\beta g)}=\alpha\widehat f+\beta\widehat g
 \qquad\text{(线性性)};
\end{equation}
\begin{equation}
\label{eq:ch1-fourier-l1-linfty}
 \|\widehat f\|_\infty\le\|f\|_1
 \quad\text{且 }\widehat f\text{ 连续};
\end{equation}
\begin{equation}
\label{eq:ch1-fourier-riemann-lebesgue}
 \lim_{|\xi|\to\infty}\widehat f(\xi)=0
 \qquad\text{(Riemann--Lebesgue)};
\end{equation}
\begin{equation}
\label{eq:ch1-fourier-convolution}
 \widehat{(f*g)}=\widehat f\,\widehat g;
\end{equation}

\hypertarget{chapter-01-page-024}{}
\begin{equation}
\label{eq:ch1-fourier-translation-modulation}
 \widehat{(\tau_hf)}(\xi)=\widehat f(\xi)e^{2\pi ih\cdot\xi},
 \quad\text{其中 }\tau_hf(x)=f(x+h);
 \qquad
 \widehat{(fe^{2\pi ih\cdot x})}(\xi)=\widehat f(\xi-h);
\end{equation}
\begin{equation}
\label{eq:ch1-fourier-orthogonal-transform}
 \text{若 }\rho\in O_n\text{(正交变换), 则 }
 \widehat{f(\rho\cdot)}(\xi)=\widehat f(\rho\xi);
\end{equation}
\begin{equation}
\label{eq:ch1-fourier-dilation}
 \text{若 }g(x)=\lambda^{-n}f(\lambda^{-1}x),
 \text{ 则 }\widehat g(\xi)=\widehat f(\lambda\xi);
\end{equation}
\begin{equation}
\label{eq:ch1-fourier-derivative}
 \widehat{\left(\frac{\partial f}{\partial x_j}\right)}(\xi)
 =2\pi i\xi_j\widehat f(\xi);
\end{equation}
\begin{equation}
\label{eq:ch1-fourier-coordinate-multiplication}
 \widehat{(-2\pi ix_jf)}(\xi)=\frac{\partial\widehat f}{\partial\xi_j}(\xi).
\end{equation}

$\widehat f$ 的连续性由控制收敛定理得到; \eqref{eq:ch1-fourier-riemann-lebesgue} 可仿照引理~\ref{lem:ch1-riemann-lebesgue} 证明; 其余性质分别由变量代换、Fubini 定理和分部积分得到. 在 \eqref{eq:ch1-fourier-derivative} 中假设 $\partial f/\partial x_j\in L^1$, 在 \eqref{eq:ch1-fourier-coordinate-multiplication} 中假设 $x_jf\in L^1$.

与环面上的情形不同, $L^1(\mathbb R^n)$ 不包含 $L^p(\mathbb R^n)$($p>1$), 因此 \eqref{eq:ch1-fourier-transform-definition} 不能定义这些空间中函数的 Fourier 变换. 出于同一原因, 本应由 $\widehat f$ 恢复 $f$ 的公式
\[
 \int_{\mathbb R^n}\widehat f(\xi)e^{2\pi ix\cdot\xi}\,d\xi
\]
可能没有意义, 因为关于 $\widehat f$ 我们只有 \eqref{eq:ch1-fourier-l1-linfty} 和 \eqref{eq:ch1-fourier-riemann-lebesgue}, 而它们并不蕴含 $\widehat f$ 可积(事实上, $\widehat f$ 一般不可积).

